\documentclass[../../main.tex]{subfiles}
\begin{document}
\chapter{Writing and organizing a chapter}
\label{ch:writing}

\begin{objectives}
\begin{itemize}
  \item Add a chapter and its references.
  \item Use the margin for supporting material.
  \item Choose a suitable teaching environment.
\end{itemize}
\end{objectives}

\section{The project structure}
The main document orders the chapters. Each chapter uses \texttt{subfiles}
to inherit the preamble. Keep packages beside \texttt{main.tex} and your
own notation in \texttt{macros/macros.tex}.

\begin{note}
From \texttt{book-template}, run
\texttt{latexmk -pdf -outdir=build main.tex}. This also updates the
bibliography with Biber.
\end{note}

To add a chapter, copy a file at the same directory depth, change its title and
labels, and add a corresponding command to \texttt{main.tex}:
\begin{lstlisting}[language={[LaTeX]TeX}]
\subfile{sections/abstract/new-chapter.tex}
\end{lstlisting}
The first line of the chapter should point to the main document:
\begin{lstlisting}[language={[LaTeX]TeX}]
\documentclass[../../main.tex]{subfiles}
\begin{document}
\chapter{Your chapter title}
% Your content goes here.
\end{document}
\end{lstlisting}

\section{Text and margin}
The main column carries the argument; the margin holds supporting notes
and figures.\sidenote{Keep a note short enough to read at a
glance. A long proof belongs in the main text.}
Use \verb|\sidenote{...}| for numbered notes and
\verb|\marginnote{...}| for unnumbered notes. Place these commands in
ordinary prose, outside theorem boxes.

The \texttt{marginfigure} and \texttt{margintable} environments hold small
illustrations and tables. Use \verb|\sidecaption{...}| inside an ordinary
figure to put its caption in the margin, as in the calculus chapter.
The \texttt{fullwidth} environment extends material across both columns;
\texttt{figure*} and \texttt{table*} do the same for floats. Leave the
adjacent margin clear when using a full-width display.

Open a chapter with a short \texttt{objectives} environment containing an
\texttt{itemize} list. It places a learning-objectives card in the margin.
The \texttt{bigidea} environment highlights a central principle, while
\texttt{recap} collects the lesson at the end. Both \texttt{bigidea} and
\texttt{exercise} accept an optional title; their numbers restart in each
chapter. For example:
\begin{lstlisting}[language={[LaTeX]TeX}]
\begin{exercise}[A useful special case]
\label{ex:special-case}
Prove the claim when n = 2.
\end{exercise}
\end{lstlisting}

\section{Definitions, examples, and references}
Use \texttt{eg} for examples and \texttt{prop} for propositions.
Definitions, theorems, lemmas, and propositions have separate counters
within each chapter; examples, corollaries, remarks, and notes are unnumbered.

\begin{definition}[Even integer]\label{def:even}
An integer $n\in\ZZ$ is \vocab{even} if there exists $k\in\ZZ$ such that $n=2k$.
\end{definition}

\begin{eg}
By \cref{def:even}, the integer $14$ is even: take $k=7$.
\end{eg}

\begin{theorem}\label{thm:even-sum}
The sum of two even integers is even.
\end{theorem}
\begin{proof}
Write $m=2a$ and $n=2b$ with $a,b\in\ZZ$. Then
$m+n=2(a+b)$, and $a+b$ is an integer.
\end{proof}

Put \verb|\label{thm:even-sum}| inside the theorem.
Reference it with \verb|\cref{thm:even-sum}| to obtain \cref{thm:even-sum}.
Use \verb|\Cref| at the start of a sentence. Label numbered objects;
an unnumbered example has no example number to reference.

\section{Citations and bibliography files}
Store reading references in \texttt{references.bib} and illustration credits
in \texttt{images.bib}. Mark credits with \texttt{keywords = \{image\}};
this keyword separates them from the main bibliography.

For example, \textcite{judson2020} is an openly available algebra text.
A parenthetical citation looks like this \parencite{concrete1994}.
Both commands use an entry key from a bibliography file:
\begin{lstlisting}[language={[LaTeX]TeX}]
\textcite{judson2020}
\parencite{concrete1994}
\end{lstlisting}

\begin{remark}
Bibliography entries appear only when cited. Use \verb|\nocite{entry-key}|
to include an uncited reading-list item before printing the bibliography.
After editing a \texttt{.bib} file, rebuild with
\texttt{latexmk}; a single LaTeX pass is not enough.
\end{remark}

\section{Adapting the design}
Edit \texttt{book-design.tex} to change the page layout, fonts, headings,
or teaching callouts. Edit \texttt{preamble.tex} to change theorem numbering
or code formatting, and \texttt{main.tex} for the title, subtitle, author,
and date.

Set \verb|\booktheme| in the preamble to \texttt{light} for Catppuccin Latte
or \texttt{dark} for Macchiato. Alternatively, build \texttt{main-light.tex}
or \texttt{main-dark.tex}; both use the same chapters and metadata.

The \texttt{macro} package supplies shortcuts such as $\RR$, $\ZZ$,
$\abs{x}$, $\norm{v}$, $\floor{x}$, and $\Hom(G,H)$. Use these commands
in math mode. The appendix collects a few more useful examples.
\end{document}
